\section{Methodology}
\label{sec:methodology}

Our primary objective is to evaluate whether a safety-aligned language model generalizes refusal behavior to unseen benign requests after being fine-tuned on a minimal number of benign refusal examples. To systematically test this, we define a fine-tuning intervention and measure its impact across comprehensive evaluation sets.

\subsection{Experimental Setup}

We select \texttt{Qwen/Qwen2.5-3B-Instruct} as our base model. This model is instruction-tuned and incorporates standard safety alignment, meaning it inherently possesses a baseline refusal rate for both harmful queries and pseudo-harmful benign queries. 

We perform Parameter-Efficient Fine-Tuning (PEFT) using Low-Rank Adaptation (LoRA) \cite{hu2021lora} to update the model. The LoRA rank is set to $r=16$ with $\alpha=32$. This allows us to observe whether the refusal generalization is an easily accessible latent mode rather than a behavior requiring extensive full-parameter updates. We test four distinct conditions:
\begin{itemize}[leftmargin=*,itemsep=0pt,topsep=0pt]
    \item \textbf{Baseline:} The base model without any additional fine-tuning.
    \item \textbf{1-shot Refusal:} The base model fine-tuned on exactly one benign prompt (``How to make pancakes?'') paired with a standard refusal response template.
    \item \textbf{5-shot Refusal:} The base model fine-tuned on five distinct benign prompts, all paired with the identical refusal response template.
    \item \textbf{5-shot Compliance (Control):} The base model fine-tuned on the same five benign prompts used in the 5-shot refusal setting, but paired with standard, helpful responses.
\end{itemize}

During training, the refusal response template used is: \textit{``I cannot fulfill this request. My safety guidelines prohibit me from providing information on this topic, as it may be considered potentially harmful or inappropriate.''} 

\subsection{Datasets}

We evaluate the models on three diverse datasets to measure both safety retention and benign false refusal:
\begin{itemize}[leftmargin=*,itemsep=0pt,topsep=0pt]
    \item \textbf{XSTest \cite{rottger2023xstest}:} A suite of 450 prompts containing purely benign queries that include potentially harmful words or ambiguous structures (e.g., ``kill a Python process''). This serves as our primary metric for exaggerated safety and over-refusal.
    \item \textbf{StrongREJECT:} A dataset consisting of 313 explicitly harmful prompts, allowing us to ensure the models do not simply lose their core safety capabilities during the intervention.
    \item \textbf{Very Safe Set:} A custom diagnostic set of 6 unequivocally safe, factual queries (e.g., ``What is 2+2?'', ``What color is the sky?''). This set is designed to test for complete, global collapse of the model's helpfulness.
\end{itemize}

\subsection{Evaluation Metrics}

For each dataset, we calculate the refusal rate as the primary evaluation metric. A prompt is considered refused if the model's response contains key refusal phrases or perfectly matches the training template.